% ── Κεφάλαιο 2: Υπόβαθρο ────────────────────────────────────

\chapter{Υπόβαθρο}
\label{ch:background}

Το κεφάλαιο αυτό παρουσιάζει το θεωρητικό υπόβαθρο που απαιτείται για την κατανόηση της προτεινόμενης αρχιτεκτονικής: από τα βασικά της μηχανικής μάθησης και των συνελικτικών δικτύων, μέχρι το Federated Learning, την αναγνώριση δραστηριότητας, τα Ψηφιακά Δίδυμα και τη θεωρία παιγνίων. Κάθε ενότητα εισάγει τις έννοιες που χρησιμοποιούνται άμεσα στην υλοποίηση και τα πειράματα των επόμενων κεφαλαίων.

% ── §2.1: Μηχανική Μάθηση και Ανάλυση Σημάτων ───────────────

\section{Μηχανική Μάθηση και Ανάλυση Σημάτων}
\label{sec:ml_signals}

Η σύγχρονη ανάλυση σημάτων αισθητήρων βασίζεται σε μεθόδους βαθιάς μάθησης
(deep learning), οι οποίες μαθαίνουν αυτόματα ιεραρχικές αναπαραστάσεις
από τα ακατέργαστα δεδομένα, χωρίς να απαιτείται χειροκίνητη μηχανική
χαρακτηριστικών (feature engineering) \cite{Goodfellow2016}. Στο πλαίσιο
της διπλωματικής αυτής, ενδιαφερόμαστε ειδικά για την ανάλυση
\textit{πολυκαναλικών χρονοσειρών} (multivariate time-series) από
επιταχυνσιόμετρα και γυροσκόπια — δεδομένα δηλαδή με έντονη τοπική
χρονική δομή και χαμηλή εξάρτηση μεγάλης εμβέλειας.

\subsection{Συνελικτικά Νευρωνικά Δίκτυα για Χρονοσειρές}
\label{subsec:cnn_timeseries}

Τα Συνελικτικά Νευρωνικά Δίκτυα (Convolutional Neural Networks, CNN)
αρχικά αναπτύχθηκαν για την ανάλυση εικόνων, αξιοποιώντας δισδιάστατα
φίλτρα (2D kernels) για την εξαγωγή χωρικών χαρακτηριστικών \cite{He2016}.
Ωστόσο, όταν τα δεδομένα εισόδου είναι χρονοσειρές — δηλαδή ακολουθίες
μετρήσεων κατά τον χρονικό άξονα — η εφαρμογή μονοδιάστατης συνέλιξης
(1D convolution) αποδεικνύεται αποτελεσματικότερη και υπολογιστικά πιο
αποδοτική \cite{Chen2021}.

Ένα \textbf{1D-CNN επίπεδο} εφαρμόζει ένα σύνολο από $F$ φίλτρα μήκους
$k$ κατά μήκος της χρονικής διάστασης. Για είσοδο $\mathbf{X} \in
\mathbb{R}^{C \times T}$ (όπου $C$ είναι ο αριθμός καναλιών και $T$ τα
χρονικά βήματα), η έξοδος του φίλτρου $f$ στη θέση $t$ υπολογίζεται ως:

\begin{equation}
    \label{eq:conv1d}
    h_f(t) = \sigma\!\left(\sum_{c=1}^{C} \sum_{\tau=0}^{k-1}
    W_{f,c,\tau} \cdot X_{c,\, t+\tau} + b_f\right)
\end{equation}

\noindent όπου $W_{f,c,\tau}$ τα εκπαιδευόμενα βάρη του φίλτρου $f$
για κανάλι $c$ και χρονική θέση $\tau$, $b_f$ η μεροληψία (bias) και
$\sigma(\cdot)$ η συνάρτηση ενεργοποίησης ReLU: $\sigma(x) = \max(0, x)$.

Κρίσιμο πλεονέκτημα του 1D-CNN έναντι εναλλακτικών αρχιτεκτονικών,
όπως τα δίκτυα LSTM \cite{Hochreiter1997}, είναι η \textit{κοινή χρήση
παραμέτρων} (parameter sharing) κατά τον χρονικό άξονα: το ίδιο φίλτρο
εφαρμόζεται σε κάθε χρονική θέση, μειώνοντας δραστικά τον αριθμό
παραμέτρων και επιταχύνοντας την εκπαίδευση. Επιπλέον, τα 1D-CNNs
εκμεταλλεύονται φυσικά την τοπική δομή των σημάτων κίνησης (π.χ.\ μια
βηματική κύκλος διαρκεί ~0.5–1 δευτερόλεπτα, κάτι που αποτυπώνεται
σε ένα παράθυρο μερικών δεκάδων δειγμάτων).

\subsection{Δεξαμενή Χαρακτηριστικών και Πλήρης Σύνδεση}
\label{subsec:pooling_fc}

Μετά από κάθε συνελικτικό επίπεδο, εφαρμόζεται \textbf{Max Pooling}
κατά μήκος της χρονικής διάστασης, το οποίο μειώνει την ανάλυση της
χρονικής αλληλουχίας κατά παράγοντα 2 (ή μεγαλύτερο), διατηρώντας
μόνο την ισχυρότερη ενεργοποίηση σε κάθε παράθυρο. Αυτό προσφέρει
\textit{μετατοπιστική αναλλοίωτη} (translation invariance): το δίκτυο
αναγνωρίζει τοπικά μοτίβα κίνησης ανεξάρτητα από την ακριβή χρονική
θέση τους μέσα στο παράθυρο.

Στο τελικό στάδιο, ο τρισδιάστατος χάρτης χαρακτηριστικών (feature map)
ισοπεδώνεται (flatten) σε διάνυσμα και οδηγείται σε πλήρως συνδεδεμένα
(fully-connected) επίπεδα, τα οποία παράγουν λογιτικές βαθμολογίες
ανά κλάση δραστηριότητας. Η τελική ταξινόμηση γίνεται μέσω \textit{softmax}:

\begin{equation}
    \label{eq:softmax}
    p_j = \frac{e^{z_j}}{\sum_{l=1}^{K} e^{z_l}}, \quad j = 1,\ldots,K
\end{equation}

\noindent όπου $z_j$ οι λογιτικές βαθμολογίες και $K = 6$ ο αριθμός
κλάσεων δραστηριότητας.

\subsection{Σύγκριση με Εναλλακτικές Αρχιτεκτονικές}
\label{subsec:architecture_comparison}

Η ολοκληρωμένη ανασκόπηση των Chen et al.\ \cite{Chen2021} αξιολογεί
συστηματικά CNN, LSTM και υβριδικές αρχιτεκτονικές για HAR. Τα
ευρήματά τους δείχνουν ότι τα 1D-CNN επιτυγχάνουν ανταγωνιστική ακρίβεια
συγκριτικά με τα LSTM, ενώ παρουσιάζουν σημαντικά μικρότερο υπολογιστικό
κόστος — ιδιότητα κρίσιμη σε FL περιβάλλοντα όπου η τοπική εκπαίδευση
γίνεται σε περιορισμένους πόρους (smartphones). Πιο πρόσφατα, μηχανισμοί
\textit{self-attention} (Transformers) \cite{Vaswani2017} έχουν εφαρμοστεί
και σε HAR χρονοσειρές, επιτυγχάνοντας υψηλή ακρίβεια, ωστόσο ο
σημαντικά μεγαλύτερος αριθμός παραμέτρων τους τα καθιστά λιγότερο
κατάλληλα για τοπική εκπαίδευση σε περιβάλλον FL με περιορισμένους πόρους.
Για τον λόγο αυτό, η παρούσα εργασία υιοθετεί αρχιτεκτονική 1D-CNN, η
οποία περιγράφεται αναλυτικά στην Ενότητα~\ref{sec:cnn_design}.

% ── §2.2: Ομοσπονδιακή Μάθηση (Federated Learning) ──────────

\section{Ομοσπονδιακή Μάθηση (Federated Learning)}
\label{sec:federated_learning}

Η Ομοσπονδιακή Μάθηση (Federated Learning, FL) είναι ένα παράδειγμα
κατανεμημένης μηχανικής μάθησης στο οποίο ένα καθολικό μοντέλο
εκπαιδεύεται από ένα πλήθος $n$ συμμετεχόντων (clients) χωρίς τα
ακατέργαστα δεδομένα τους να εγκαταλείπουν ποτέ τη συσκευή τους
\cite{McMahan2017}. Αντί για δεδομένα, μεταδίδονται μόνο ενημερώσεις
μοντέλου (model updates) στον κεντρικό συντονιστή (server), ο οποίος
τις συνθέτει σε ένα βελτιωμένο καθολικό μοντέλο.

\subsection{Αρχιτεκτονική Server-Client}
\label{subsec:fl_architecture}

Το τυπικό FL σύστημα αποτελείται από έναν \textbf{server} και $n$
\textbf{clients}. Ο server αρχικοποιεί ένα καθολικό μοντέλο
$\mathbf{w}^{(0)}$ και ορίζει τη στρατηγική επιλογής clients ανά γύρο
επικοινωνίας. Κάθε client $i$ διατηρεί τοπικό σύνολο δεδομένων
$\mathcal{D}_i$ με $n_i = |\mathcal{D}_i|$ δείγματα. Η κατανεμημένη
βελτιστοποίηση στοχεύει στην ελαχιστοποίηση της σταθμισμένης συνολικής
απώλειας:

\begin{equation}
    \label{eq:fl_objective}
    \min_{\mathbf{w}} \; F(\mathbf{w}) = \sum_{i=1}^{n}
    \frac{n_i}{N} \, F_i(\mathbf{w}), \quad N = \sum_{i=1}^{n} n_i
\end{equation}

\noindent όπου $F_i(\mathbf{w}) = \frac{1}{n_i}\sum_{(\mathbf{x},y)
\in \mathcal{D}_i} \ell(\mathbf{w}; \mathbf{x}, y)$ είναι η τοπική
απώλεια του client $i$ και $\ell(\cdot)$ η συνάρτηση κόστους
(cross-entropy για ταξινόμηση).

\subsection{Ο Αλγόριθμος FedAvg}
\label{subsec:fedavg}

Ο αλγόριθμος Federated Averaging (FedAvg), που εισήχθη από τους McMahan
et al.\ \cite{McMahan2017}, αποτελεί τον de facto αλγόριθμο αναφοράς
του FL. Επαναλαμβάνει τα ακόλουθα βήματα για $R$ γύρους επικοινωνίας:

\begin{enumerate}
    \item \textbf{Κοινοποίηση:} Ο server αποστέλλει τις τρέχουσες
    παραμέτρους $\mathbf{w}^{(r)}$ στο υποσύνολο $\mathcal{S}^{(r)}$
    επιλεγμένων clients (με $|\mathcal{S}^{(r)}| = m \leq n$).

    \item \textbf{Τοπική εκπαίδευση:} Κάθε client $i \in \mathcal{S}^{(r)}$
    εκτελεί $E$ εποχές τοπικής βελτιστοποίησης (SGD ή Adam) στα
    δεδομένα του, παράγοντας ενημερωμένες παραμέτρους
    $\mathbf{w}_i^{(r+1)}$.

    \item \textbf{Συνάθροιση:} Ο server συνθέτει τις ενημερώσεις
    στάθμισε με βάση το μέγεθος των τοπικών συνόλων:
    \begin{equation}
        \label{eq:fedavg}
        \mathbf{w}^{(r+1)} = \sum_{i \in \mathcal{S}^{(r)}}
        \frac{n_i}{\sum_{j \in \mathcal{S}^{(r)}} n_j}
        \mathbf{w}_i^{(r+1)}
    \end{equation}
\end{enumerate}

Η θεωρητική ανάλυση σύγκλισης του FedAvg υπό Non-IID κατανομές
παρουσιάζεται από τους Li et al.\ \cite{Li2020}, οι οποίοι αποδεικνύουν
ότι η ετερογένεια δεδομένων μπορεί να επιβραδύνει σημαντικά τη σύγκλιση
λόγω του φαινομένου \textit{client drift}. Το Σχήμα~\ref{fig:fl_round}
απεικονίζει σχηματικά τον κύκλο ενός γύρου FL.

\begin{figure}[H]
\centering
\begin{tikzpicture}[
    node distance=1.5cm and 2.2cm,
    box/.style={rectangle, rounded corners=4pt, draw=black, thick,
                minimum width=2.8cm, minimum height=0.9cm,
                align=center, font=\small},
    server/.style={box, fill=blue!15},
    client/.style={box, fill=orange!20},
    arrow/.style={-{Stealth[length=6pt]}, thick},
    dasharrow/.style={-{Stealth[length=6pt]}, thick, dashed}
]

% Server
\node[server] (server) {FL Server\\[2pt]$\mathbf{w}^{(r)}$};

% Clients
\node[client, below left=1.8cm and 2.5cm of server]  (c1) {Client 1\\[2pt]$\mathcal{D}_1$};
\node[client, below=1.8cm of server]                  (c2) {Client 2\\[2pt]$\mathcal{D}_2$};
\node[client, below right=1.8cm and 2.5cm of server]  (c3) {Client $N$\\[2pt]$\mathcal{D}_N$};

% Dots between clients
\node[font=\Large] at ($(c2)!0.5!(c3) + (0.6,0)$) {$\cdots$};

% Broadcast arrows (server → clients)
\draw[arrow, blue!60] (server) -- node[left,  font=\scriptsize]{$\mathbf{w}^{(r)}$} (c1);
\draw[arrow, blue!60] (server) -- node[right, font=\scriptsize]{$\mathbf{w}^{(r)}$} (c2);
\draw[arrow, blue!60] (server) -- node[right, font=\scriptsize]{$\mathbf{w}^{(r)}$} (c3);

% Local training labels
\node[font=\scriptsize, text=gray] at ($(c1)+(0,-0.7)$) {$E$ τοπικές εποχές};
\node[font=\scriptsize, text=gray] at ($(c2)+(0,-0.7)$) {$E$ τοπικές εποχές};
\node[font=\scriptsize, text=gray] at ($(c3)+(0,-0.7)$) {$E$ τοπικές εποχές};

% Aggregation node
\node[server, below=3.8cm of server] (agg) {FedAvg\\[2pt]Aggregation};

% Upload arrows (clients → aggregation)
\draw[arrow, orange!70!black] (c1) -- node[left,  font=\scriptsize]{$\mathbf{w}_1^{(r+1)}$} (agg);
\draw[arrow, orange!70!black] (c2) -- node[right, font=\scriptsize]{$\mathbf{w}_2^{(r+1)}$} (agg);
\draw[arrow, orange!70!black] (c3) -- node[right, font=\scriptsize]{$\mathbf{w}_N^{(r+1)}$} (agg);

% Loop arrow
\draw[dasharrow, green!50!black, bend right=50]
      (agg.east) to node[right, align=left, font=\scriptsize]{$\mathbf{w}^{(r+1)}$\\[1pt]επόμενος γύρος} (server.east);

\end{tikzpicture}
\caption{Κύκλος ενός FL γύρου (FedAvg): ο server εκπέμπει τις τρέχουσες παραμέτρους, κάθε client εκτελεί $E$ τοπικές εποχές χωρίς κοινοποίηση δεδομένων, και τα ενημερωμένα βάρη συνδυάζονται σταθμισμένα για το επόμενο global μοντέλο.}
\label{fig:fl_round}
\end{figure}

\subsection{Προκλήσεις του Federated Learning}
\label{subsec:fl_challenges}

Η ολοκληρωμένη ανασκόπηση των Kairouz et al.\ \cite{Kairouz2021}
κατατάσσει τις κύριες προκλήσεις του FL σε τρεις κατηγορίες:

\begin{itemize}
    \item \textbf{Ετερογένεια δεδομένων (Non-IID):} Τα δεδομένα κάθε
    client ακολουθούν διαφορετική κατανομή, καθώς αντικατοπτρίζουν
    τις ατομικές συνήθειες και κινητικά μοτίβα κάθε χρήστη. Αυτό
    αποκλίνει από τη βέλτιστη υπόθεση IID (Independent and Identically
    Distributed) της στατιστικής θεωρίας \cite{Li2020}.

    \item \textbf{Κόστος επικοινωνίας:} Κάθε γύρος απαιτεί μετάδοση
    ολόκληρου του μοντέλου (ή των παραμέτρων του) μεταξύ server και
    clients. Για μεγάλα μοντέλα ή αργές συνδέσεις, αυτό αποτελεί
    σημαντικό περιορισμό \cite{Konecny2016}.

    \item \textbf{Αξιοπιστία clients:} Clients μπορεί να αποσυνδεθούν
    κατά τη διάρκεια ενός γύρου (stragglers) ή να υποβάλουν ψευδείς
    ενημερώσεις (Byzantine clients). Η παρούσα εργασία αντιμετωπίζει
    την πρώτη κατηγορία μέσω του μηχανισμού κινήτρων
    (Κεφάλαιο~\ref{ch:implementation}).
\end{itemize}

\subsection{Ιδιωτικότητα στο FL}
\label{subsec:fl_privacy}

Παρόλο που το FL δεν μεταδίδει ακατέργαστα δεδομένα, οι ενημερώσεις
μοντέλου μπορούν υπό συνθήκες να αποκαλύψουν ευαίσθητες πληροφορίες
μέσω \textit{gradient inversion attacks} \cite{Shokri2015}. Μέθοδοι
όπως η Differential Privacy \cite{Abadi2016} και η ασφαλής συνάθροιση
\cite{Bonawitz2017} αμβλύνουν αυτόν τον κίνδυνο, αν και βρίσκονται
εκτός του άμεσου πεδίου της παρούσας εργασίας.

% ── §2.3: Αναγνώριση Δραστηριότητας Ανθρώπου (HAR) ──────────

\section{Αναγνώριση Δραστηριότητας Ανθρώπου (HAR)}
\label{sec:har}

Η Αναγνώριση Δραστηριότητας Ανθρώπου (Human Activity Recognition, HAR)
αφορά την αυτόματη κατάταξη της φυσικής δραστηριότητας ενός ατόμου σε
προκαθορισμένες κλάσεις, βάσει δεδομένων αισθητήρων κινητής συσκευής
\cite{Anguita2013}. Στο πλαίσιο των συστημάτων AAL, το HAR επιτρέπει
την παρακολούθηση της λειτουργικής αυτονομίας χρηστών σε πραγματικό
χρόνο, χωρίς να απαιτείται φορετός εξοπλισμός ή υποδομή αισθητήρων
στο περιβάλλον \cite{Liu2021}.

\subsection{Αισθητήρες Κινητών Συσκευών}
\label{subsec:sensors}

Τα σύγχρονα smartphones διαθέτουν ενσωματωμένους αισθητήρες υψηλής
ευαισθησίας που καταγράφουν κινηματικές πληροφορίες σε πραγματικό χρόνο:

\begin{itemize}
    \item \textbf{Επιταχυνσιόμετρο (Accelerometer):} Μετρά την γραμμική
    επιτάχυνση κατά τους τρεις άξονες $\{x, y, z\}$ σε $\text{m/s}^2$.
    Συλλαμβάνει δυναμικά μοτίβα (βηματισμός, πτώση) αλλά και στατικές
    κατευθύνσεις (βαρύτητα ανάλογα με κλίση συσκευής).

    \item \textbf{Γυροσκόπιο (Gyroscope):} Μετρά τον ρυθμό γωνιακής
    στροφής κατά τους τρεις άξονες σε $\text{rad/s}$. Συμπληρώνει
    το επιταχυνσιόμετρο αποτυπώνοντας περιστροφικές κινήσεις που δεν
    ανιχνεύονται από τη γραμμική επιτάχυνση.
\end{itemize}

Ο συνδυασμός των δύο αισθητήρων παράγει ένα σήμα 6 καναλιών. Στο
UCI HAR dataset \cite{Anguita2013}, προστίθενται 3 ακόμα κανάλια
(επιταχύνσεις σώματος μετά από αφαίρεση βαρύτητας), οδηγώντας σε
είσοδο 9 καναλιών για τα αντίστοιχα μοντέλα μάθησης.

\subsection{Τεμαχισμός Σήματος με Συρόμενο Παράθυρο}
\label{subsec:sliding_window_bg}

Η ακατέργαστη χρονοσειρά τεμαχίζεται σε παράθυρα σταθερού μήκους με
τη μέθοδο \textit{sliding window}. Για μια χρονοσειρά $\mathbf{s}(t)$
δειγματοληψίας $f_s = 50\,\text{Hz}$, το παράθυρο $\mathbf{X}^{(k)}$
αρχίζει στο δείγμα $k \cdot (1-\rho) \cdot W$ και λήγει $W$ δείγματα
αργότερα, όπου $W$ το μήκος παραθύρου (σε δείγματα) και $\rho$ η
αναλογία επικάλυψης (overlap). Στο UCI HAR, $W = 128$ δείγματα ($2.56\,\text{s}$)
και $\rho = 0.5$, αποδίδοντας $\lfloor 2(T/W - 1) \rfloor + 1$ παράθυρα
ανά χρονοσειρά.

Κάθε παράθυρο $\mathbf{X}^{(k)} \in \mathbb{R}^{9 \times 128}$
αντιστοιχεί σε ένα δείγμα εισόδου για τον ταξινομητή. Αυτή ακριβώς
είναι η μορφή εισόδου $[B \times 9 \times 128]$ που δέχεται η
αρχιτεκτονική 1D-CNN (βλ.\ Ενότητα~\ref{sec:cnn_design}).

\subsection{Κλάσεις Δραστηριότητας}
\label{subsec:activity_classes}

Το UCI HAR dataset \cite{Anguita2013} περιλαμβάνει δεδομένα από 30
εθελοντές (ηλικία 19--48 ετών) για $K = 6$ κλάσεις δραστηριότητας:
\textit{περπάτημα} (Walking), \textit{ανάβαση σκαλιών} (Walking Upstairs),
\textit{κατάβαση σκαλιών} (Walking Downstairs), \textit{κάθισμα}
(Sitting), \textit{όρθια στάση} (Standing) και \textit{ξάπλωμα}
(Laying). Η κατηγοριοποίηση αυτή — τρεις δυναμικές και τρεις στατικές
κλάσεις — αποτελεί ευρέως αποδεκτό πρότυπο αξιολόγησης στη βιβλιογραφία
HAR \cite{Chen2021}.

\subsection{Προκλήσεις Αναγνώρισης σε Πραγματικές Συνθήκες}
\label{subsec:har_challenges}

Παρά τις επιδόσεις σε ελεγχόμενα πειραματικά περιβάλλοντα, το HAR
υπό πραγματικές συνθήκες AAL αντιμετωπίζει σημαντικές προκλήσεις
\cite{Hammerla2016}:

\begin{itemize}
    \item \textbf{Ατομική μεταβλητότητα:} Διαφορετικοί χρήστες
    εκτελούν ίδιες δραστηριότητες με εντελώς διαφορετικά κινηματικά
    μοτίβα — λόγω σωματότυπου, ηλικίας ή τρόπου κρατήματος της
    συσκευής. Αυτό οδηγεί σε Non-IID κατανομές δεδομένων μεταξύ
    clients στο FL σύστημα.

    \item \textbf{Θόρυβος αισθητήρων:} Ο θόρυβος αυξάνεται σε
    πραγματικά AAL σενάρια (άτυπη χρήση τηλεφώνου, κίνηση στο
    περιβάλλον) σε αντίθεση με τα πρωτόκολλα ελεγχόμενης συλλογής.

    \item \textbf{Ανισορροπία κλάσεων:} Στατικές δραστηριότητες
    (κάθισμα, ξάπλωμα) καταλαμβάνουν δυσανάλογα μεγαλύτερο χρόνο
    για ηλικιωμένους χρήστες, δημιουργώντας ανισορροπία στα τοπικά
    datasets των clients.
\end{itemize}

Οι ανωτέρω προκλήσεις επιβεβαιώνουν την ανάγκη για αξιολόγηση σε
πραγματικά δεδομένα — όπως το AAL dataset \cite{AAL_UCI2016}
— πέρα από τα ελεγχόμενα benchmarks. Ένας πρόσθετος παράγοντας
που επηρεάζει ειδικά τα FL συστήματα HAR είναι η ανάπτυξη σε
συσκευές περιορισμένων πόρων (smartphones, wearables): η υπολογιστική
ισχύς και η μνήμη είναι περιορισμένες, απαιτώντας αρχιτεκτονικές
μοντέλων που ισορροπούν ακρίβεια και αποδοτικότητα. Η εργασία των
Yao et al.\ \cite{Yao2017} (DeepIoT) δείχνει ότι η συμπίεση μοντέλων
για IoT συσκευές είναι εφικτή χωρίς σημαντική απώλεια ακρίβειας —
μια κατεύθυνση συμπληρωματική με αυτή της παρούσας εργασίας, η
οποία επιλέγει ήδη αρχιτεκτονική με περιορισμένο αριθμό παραμέτρων
($\approx 47$k) για τοπική εκπαίδευση. Παράλληλα, η χρήση πολλαπλών
τύπων αισθητήρων (fusion) για βελτίωση της αναγνώρισης σε ρεαλιστικά
σενάρια εξετάζεται από τους Gu et al.\ \cite{Gu2022}, αναδεικνύοντας
ότι η επιλογή των καναλιών εισόδου επηρεάζει σημαντικά την απόδοση
σε Non-IID κατανομές.

% ── §2.4: Ψηφιακός Δίδυμος Χρήστη (Digital Twin) ────────────

\section{Ψηφιακός Δίδυμος Χρήστη (Digital Twin)}
\label{sec:digital_twin}

Η έννοια του Ψηφιακού Διδύμου (Digital Twin, DT) ορίζεται ως μια
δυναμική ψηφιακή αναπαράσταση ενός φυσικού οντότητας ή συστήματος,
η οποία ενημερώνεται συνεχώς από ροές δεδομένων πραγματικού χρόνου
\cite{Grieves2014}. Αρχικά αναπτύχθηκε στον χώρο της βιομηχανικής
παραγωγής, αλλά η εφαρμογή του στη ψηφιακή υγεία και τα AAL
συστήματα αναδεικνύεται ως πεδίο υψηλής ερευνητικής δραστηριότητας
\cite{Tao2019}.

\subsection{Στρώματα Αρχιτεκτονικής}
\label{subsec:dt_layers}

Ένας Ψηφιακός Δίδυμος χρήστη σε AAL περιβάλλον οργανώνεται σε τρία
εννοιολογικά στρώματα \cite{LiuDT2021}:

\begin{enumerate}
    \item \textbf{Στρώμα Δεδομένων (Data Layer):} Συλλέγει ακατέργαστα
    δεδομένα αισθητήρων από το smartphone (επιταχυνσιόμετρο, γυροσκόπιο)
    και άλλες πηγές (wearables, έξυπνο σπίτι). Αποτελεί την \textit{«αίσθηση»}
    του διδύμου — το τι συμβαίνει στον φυσικό κόσμο.

    \item \textbf{Στρώμα Μοντέλου (Model Layer):} Εκτελεί επεξεργασία
    και ταξινόμηση των δεδομένων μέσω μοντέλων μηχανικής μάθησης.
    Στην παρούσα εργασία, το στρώμα αυτό υλοποιείται από το 1D-CNN
    μοντέλο HAR που εκπαιδεύεται συνεργατικά μέσω FL.

    \item \textbf{Στρώμα Υπηρεσιών (Service Layer):} Παρέχει ανθρώπινα
    αναγνώσιμες πληροφορίες — ειδοποιήσεις, τάσεις δραστηριότητας,
    αναφορές σε φροντιστές — βάσει των εξόδων του στρώματος μοντέλου.
\end{enumerate}

\subsection{Ψηφιακός Δίδυμος στο FL Σύστημα}
\label{subsec:dt_in_fl}

Στο πλαίσιο της παρούσας εργασίας, κάθε client του FL συστήματος
αντιστοιχεί σε έναν χρήστη του AAL περιβάλλοντος. Το τοπικό μοντέλο
HAR κάθε client λειτουργεί ως \textit{εξατομικευμένος Ψηφιακός Δίδυμος}:
αποτυπώνει τα ιδιαίτερα κινηματικά μοτίβα του συγκεκριμένου χρήστη
χωρίς να μεταδίδει ευαίσθητα δεδομένα εκτός της συσκευής.

Η έννοια του Ψηφιακού Διδύμου επιβεβαιώνει τη σχεδιαστική επιλογή
να διατηρούμε τοπικά δεδομένα στις συσκευές: ο DT κάθε χρήστη
\textit{ζει} στη συσκευή του, ενώ η γνώση (μοντέλο) μεταδίδεται στον
server μέσω FL. Αυτή η αρχιτεκτονική εκμεταλλεύεται τα πλεονεκτήματα
τόσο του FL (προστασία ιδιωτικότητας) όσο και του DT (εξατομίκευση
και συνεχής ενημέρωση). Η συνολική αρχιτεκτονική DT+FL αποτυπώνεται
στο Σχήμα~\ref{fig:dt_fl_arch}.

\begin{figure}[H]
\centering
\begin{tikzpicture}[
    node distance=1cm and 1.8cm,
    sensor/.style={rectangle, rounded corners=3pt, draw=gray!70, thick,
                   fill=gray!10, minimum width=2cm, minimum height=0.7cm,
                   font=\scriptsize, align=center},
    dtbox/.style={rectangle, rounded corners=5pt, draw=blue!60, thick,
                  fill=blue!8, minimum width=3.2cm, minimum height=0.8cm,
                  font=\small, align=center},
    srvbox/.style={rectangle, rounded corners=5pt, draw=green!60!black, thick,
                   fill=green!10, minimum width=3.5cm, minimum height=0.9cm,
                   font=\small, align=center},
    arrow/.style={-{Stealth[length=6pt]}, thick},
    dbarrow/.style={<->, thick, dashed, gray!70}
]

% Physical world (left column)
\node[sensor] (phone1) {Smartphone 1\\[2pt](sensors)};
\node[sensor, below=0.5cm of phone1] (phone2) {Smartphone 2\\[2pt](sensors)};
\node[font=\scriptsize, below=0.3cm of phone2] (dots) {$\vdots$};
\node[sensor, below=0.3cm of dots] (phoneN) {Smartphone $N$\\[2pt](sensors)};

% DT nodes (middle column)
\node[dtbox, right=2.2cm of phone1] (dt1) {DT User 1\\[2pt](local 1D-CNN)};
\node[dtbox, right=2.2cm of phone2] (dt2) {DT User 2\\[2pt](local 1D-CNN)};
\node[font=\small, right=2.2cm of dots] (dots2) {$\vdots$};
\node[dtbox, right=2.2cm of phoneN] (dtN) {DT User $N$\\[2pt](local 1D-CNN)};

% Server (right)
\node[srvbox, right=2.5cm of dt2] (srv) {FL Server\\[2pt]FedAvg + Incentives};

% Sensor → DT (raw data stays local)
\draw[dbarrow] (phone1) -- node[above, font=\tiny]{raw data} (dt1);
\draw[dbarrow] (phone2) -- node[above, font=\tiny]{raw data} (dt2);
\draw[dbarrow] (phoneN) -- node[above, font=\tiny]{raw data} (dtN);

% DT → Server (only model weights)
\draw[arrow, blue!60] (dt1.east) -- node[above, font=\tiny]{$\Delta\mathbf{w}_1$} (srv.west|-dt1);
\draw[arrow, blue!60] (dt2.east) -- node[above, font=\tiny]{$\Delta\mathbf{w}_2$} (srv.west|-dt2);
\draw[arrow, blue!60] (dtN.east) -- node[above, font=\tiny]{$\Delta\mathbf{w}_N$} (srv.west|-dtN);

% Server → DT (global model)
\draw[arrow, green!60!black] (srv.west|-dt1) -- +(-0.4,0) -- (dt1.east);

% Labels
\node[font=\tiny, text=gray, below=0.15cm of dt1] {Τοπική εκπαίδευση};
\node[font=\tiny, text=red!60!black, above=0.15cm of phone1] {Δεδομένα: τοπικά};
\node[font=\tiny, text=green!60!black, above right=0.1cm and -0.3cm of srv]
    {Global model};

\end{tikzpicture}
\caption{Αρχιτεκτονική DT+FL: κάθε χρήστης AAL αντιστοιχεί σε έναν Ψηφιακό Δίδυμο που αποθηκεύει τοπικά τα ακατέργαστα δεδομένα αισθητήρων και στέλνει στον FL server μόνο ενημερώσεις βαρών μοντέλου ($\Delta\mathbf{w}_i$). Ο server συγκεντρώνει, αξιολογεί τη συνεισφορά και κατανέμει ανταμοιβές~\cite{LiuDT2021, Tao2019}.}
\label{fig:dt_fl_arch}
\end{figure}

% ── §2.5: Θεωρία Παιγνίων και Μηχανισμοί Κινήτρων ───────────

\section{Θεωρία Παιγνίων και Μηχανισμοί Κινήτρων}
\label{sec:game_theory}

Σε ένα FL σύστημα, η συμμετοχή των clients δεν είναι δεδομένη:
κάθε client επωμίζεται πραγματικό κόστος (κατανάλωση μπαταρίας,
υπολογιστικοί πόροι, χρόνος επικοινωνίας) και, απουσία κινήτρων,
έχει λόγο να αρνηθεί ή να συμμετάσχει με χαμηλή ποιότητα
(free-riding) \cite{Kang2019}. Η Θεωρία Παιγνίων παρέχει το
μαθηματικό πλαίσιο για τον σχεδιασμό μηχανισμών κινήτρων που
εξασφαλίζουν εθελοντική, ποιοτική συμμετοχή \cite{Zhan2023}.

\subsection{Το Παίγνιο Stackelberg}
\label{subsec:stackelberg}

Το \textbf{Παίγνιο Stackelberg} (Stackelberg Game) είναι ένα
παίγνιο δύο επιπέδων με ασύμμετρους ρόλους \cite{Sarikaya2020}:

\begin{itemize}
    \item \textbf{Ηγέτης (Leader):} Ο server του FL, ο οποίος
    ανακοινώνει πρώτος τη στρατηγική ανταμοιβής — δηλαδή, πόση
    αμοιβή $r_i$ θα λάβει ο client $i$ για τη συμμετοχή του.

    \item \textbf{Ακόλουθοι (Followers):} Οι clients, οι οποίοι
    παρατηρούν τη στρατηγική του server και επιλέγουν το επίπεδο
    προσπάθειας $e_i$ που μεγιστοποιεί τη δική τους ωφελιμότητα.
\end{itemize}

Η ωφελιμότητα (utility) κάθε client $i$ ορίζεται ως:

\begin{equation}
    \label{eq:client_utility}
    u_i = r_i \cdot q_i(e_i) - c_i(e_i)
\end{equation}

\noindent όπου $r_i$ η αμοιβή που ορίζει ο server, $q_i(e_i)$ η
ποιότητα της συνεισφοράς (αύξουσα συνάρτηση της προσπάθειας) και
$c_i(e_i)$ το κόστος προσπάθειας (αύξουσα κυρτή συνάρτηση). Ο
client επιλέγει $e_i^* = \arg\max_{e_i} u_i$, οδηγώντας σε
βέλτιστη απόκριση ανάλογη της αμοιβής.

\subsection{Ισορροπία Stackelberg και Βελτιστοποίηση Server}
\label{subsec:stackelberg_eq}

Γνωρίζοντας ότι οι clients θα ακολουθήσουν τη βέλτιστη απόκρισή
τους $e_i^*(r_i)$, ο server επιλύει ένα πρόβλημα βελτιστοποίησης
πρώτου επιπέδου:

\begin{equation}
    \label{eq:server_utility}
    \max_{\{r_i\}} \; U_{\text{server}} = V\!\left(\sum_{i} q_i(e_i^*(r_i))\right)
    - \sum_{i} r_i \cdot q_i(e_i^*(r_i))
\end{equation}

\noindent όπου $V(\cdot)$ η αξία της συνολικής ποιότητας συμμετοχής
για τη βελτίωση του καθολικού μοντέλου, μείον το συνολικό κόστος
αμοιβών. Η λύση του προβλήματος παρέχει τις βέλτιστες αμοιβές
$\{r_i^*\}$ που μεγιστοποιούν το κέρδος του server ενώ παράλληλα
εξασφαλίζουν την ενεργό συμμετοχή των clients.

\subsection{Εφαρμογή στο FL Σύστημα}
\label{subsec:game_theory_fl}

Στην παρούσα εργασία, ο μηχανισμός κινήτρων βασισμένος στο Παίγνιο
Stackelberg υλοποιείται στον εξυπηρετητή (server) του custom FL πλαισίου
(Κεφάλαιο~\ref{ch:implementation}). Η ποιότητα συνεισφοράς $q_i$
υπολογίζεται μέσω της αξιολόγησης Shapley (§~\ref{sec:shapley}),
δημιουργώντας έναν κλειστό βρόχο: καλύτερη συνεισφορά → υψηλότερη
αμοιβή → κίνητρο για ποιοτική συμμετοχή \cite{Pavlidis2021,Yu2020}.
Η δομή του παιγνίου αποτυπώνεται στο Σχήμα~\ref{fig:stackelberg}.

\begin{figure}[H]
\centering
\begin{tikzpicture}[
    node distance=1.2cm and 3cm,
    levelbox/.style={rectangle, rounded corners=5pt, draw=black, thick,
                     minimum width=3.5cm, minimum height=1cm,
                     align=center, font=\small},
    leader/.style={levelbox, fill=blue!15},
    follower/.style={levelbox, fill=orange!20},
    arrow/.style={-{Stealth[length=7pt]}, thick},
    annot/.style={font=\scriptsize, text=gray, align=center}
]

% Leader
\node[leader] (server) {\textbf{Server (Leader)}\\[2pt]Ανακοινώνει $\{r_i^*\}$};

% Step 1 label
\node[annot, right=0.3cm of server] {Βήμα 1: βέλτιστη\\πολιτική ανταμοιβής};

% Clients level
\node[follower, below left=2cm and 1.8cm of server]  (c1) {Client 1\\[2pt]$e_1^* = r_1/(2c)$};
\node[follower, below=2cm of server]                  (c2) {Client 2\\[2pt]$e_2^* = r_2/(2c)$};
\node[follower, below right=2cm and 1.8cm of server]  (cN) {Client $N$\\[2pt]$e_N^* = r_N/(2c)$};
\node[annot] at ($(c2)!0.5!(cN)+(0.8,0)$) {$\cdots$};

% Arrows server → clients
\draw[arrow, blue!60] (server) -- node[left,  annot]{$r_1$} (c1);
\draw[arrow, blue!60] (server) -- node[right, annot]{$r_2$} (c2);
\draw[arrow, blue!60] (server) -- node[right, annot]{$r_N$} (cN);

% Step 2 label
\node[annot] at ($(c1)+(0,-0.9)$) {Βήμα 2: βέλτιστη\\προσπάθεια $e_i^*$};

% Outcome box
\node[leader, below=3.8cm of server, minimum width=5cm] (eq)
      {\textbf{Stackelberg Ισορροπία}\\[2pt]$u_i = r_i \cdot q_i(e_i^*) - c\,(e_i^*)^2$};

% Return arrows
\draw[arrow, orange!70!black] (c1) -- node[left,  annot]{$q_1(e_1^*)$} (eq);
\draw[arrow, orange!70!black] (c2) -- node[right, annot]{$q_2(e_2^*)$} (eq);
\draw[arrow, orange!70!black] (cN) -- node[right, annot]{$q_N(e_N^*)$} (eq);

\end{tikzpicture}
\caption{Δομή του Παιγνίου Stackelberg στο FL σύστημα. Ο server (Leader) ανακοινώνει πρώτος τη βέλτιστη πολιτική ανταμοιβών $\{r_i^*\}$. Κάθε client (Follower) αποκρίνεται επιλέγοντας το επίπεδο προσπάθειας $e_i^*$ που μεγιστοποιεί τη δική του ωφελιμότητα. Η ισορροπία είναι μοναδική και αποδοτική~\cite{Sarikaya2020}.}
\label{fig:stackelberg}
\end{figure}

% ── §2.6: Αξιολόγηση Συνεισφοράς (Shapley Values) ───────────

\section{Αξιολόγηση Συνεισφοράς (Shapley Values)}
\label{sec:shapley}

Ένα κεντρικό ερώτημα στα συνεργατικά συστήματα είναι: \textit{πόσο
συνέβαλε κάθε συμμετέχων στο συνολικό αποτέλεσμα;} Στο πλαίσιο του
FL, αυτό αντιστοιχεί στο ερώτημα πόσο βελτίωσε κάθε client το
καθολικό μοντέλο. Η απάντηση δίνεται μέσω της θεωρίας Συνεργατικών
Παιγνίων (Cooperative Game Theory) και ειδικότερα μέσω των
\textit{τιμών Shapley} \cite{Shapley1953}.

\subsection{Συνεργατικά Παίγνια και Συναρτήσεις Χαρακτηριστικής}
\label{subsec:cooperative_games}

Ένα συνεργατικό παίγνιο ορίζεται ως ζεύγος $(N, v)$, όπου
$N = \{1, \ldots, n\}$ το σύνολο παικτών (clients) και
$v: 2^N \rightarrow \mathbb{R}$ η \textbf{συνάρτηση χαρακτηριστικής}
(characteristic function) που αποδίδει μια αξία σε κάθε υποσύνολο
$S \subseteq N$ \cite{Wang2019}. Στο HAR-FL σύστημά μας, $v(S)$
αντιστοιχεί στην ακρίβεια του μοντέλου που θα προέκυπτε αν μόνο
το υποσύνολο $S$ clients συμμετείχε στην εκπαίδευση.

\subsection{Ορισμός και Ιδιότητες Shapley Value}
\label{subsec:shapley_def}

Η τιμή Shapley $\phi_i(v)$ αποδίδει στον παίκτη $i$ το σταθμισμένο
μέσο περιθώριο συνεισφοράς σε όλες τις δυνατές σειρές εισόδου:

\begin{equation}
    \label{eq:shapley}
    \phi_i(v) = \sum_{S \subseteq N \setminus \{i\}}
    \frac{|S|!\,(|N|-|S|-1)!}{|N|!}
    \bigl[v(S \cup \{i\}) - v(S)\bigr]
\end{equation}

\noindent Ο παράγοντας $[v(S \cup \{i\}) - v(S)]$ μετράει την
\textit{οριακή συνεισφορά} (marginal contribution) του client $i$
όταν εντάσσεται στο υποσύνολο $S$. Ο συντελεστής
$|S|!\,(|N|-|S|-1)!/|N|!$ σταθμίζει αναλόγως με την πιθανότητα
εμφάνισης αυτής της σειράς.

Οι τιμές Shapley ικανοποιούν τέσσερις αξιωματικές ιδιότητες:
\textit{αποδοτικότητα} (το σύνολο της αξίας κατανέμεται πλήρως),
\textit{συμμετρία} (ίδια συνεισφορά → ίδια ανταμοιβή),
\textit{μηδενικός παίκτης} (μηδενική συνεισφορά → μηδενική ανταμοιβή)
και \textit{προσθετικότητα}. Αυτές οι ιδιότητες εγγυώνται τη
\textbf{δικαιοσύνη} (fairness) της κατανομής ανταμοιβών \cite{Shapley1953}.

\subsection{Υπολογιστική Πολυπλοκότητα και Προσεγγίσεις}
\label{subsec:shapley_approx}

Ο ακριβής υπολογισμός της Eq.~\eqref{eq:shapley} απαιτεί
αξιολόγηση $v(S)$ για $2^n$ υποσύνολα — εκθετική πολυπλοκότητα
που καθιστά τον ακριβή υπολογισμό απαγορευτικό για $n > 20$.
Στην πράξη χρησιμοποιείται η \textbf{προσέγγιση Leave-One-Out (LOO)}:

\begin{equation}
    \label{eq:loo}
    \hat{\phi}_i = v(N) - v(N \setminus \{i\})
\end{equation}

\noindent η οποία υπολογίζει την οριακή συνεισφορά κάθε client
ως διαφορά ακρίβειας με και χωρίς τον client $i$, απαιτώντας μόνο
$n+1$ αξιολογήσεις \cite{Gopinathan2023}. Εναλλακτικά, η
\textit{gradient-based} προσέγγιση εκτιμά τη συνεισφορά από τις
κλίσεις (gradients) χωρίς επαναξιολόγηση μοντέλου \cite{Pavlidis2021}.

Στην παρούσα εργασία, η LOO προσέγγιση χρησιμοποιείται ως η
ποσότητα $q_i$ του μηχανισμού κινήτρων (§~\ref{sec:game_theory}),
υπολογιζόμενη στον server μετά από κάθε γύρο FL
(βλ.\ Ενότητα~\ref{sec:incentive_mechanism}).

